\scanpage{4}
\mspage{2}

\unsure{Comme|Alors} $\pi$ opère sur \deleted{$\unsure{H^iX_{\bar K}}$} $E$ par ``transport de structure'', on trouve une représentation $\ell$-adique $(*)$.

\uline{Exemple} 0.4. Soit $A$ un \deleted{\unsure{schéma en}} groupe algébrique \inserted{comm.} sur $K$.
\deleted{\unsure{comm.}} alors on regarde
\[
 E=T_\ell(A(\bar K))\overset{\mathrm{dfn}}=\varprojlim_\nu {}_{\ell^\nu}A(\bar K)
\]
\deleted{\unsure{Comme}} le groupe ne change pas si on remplace $A$ par $(A^\circ)_{\unsure{\mathrm{red}}}$ (NB. $T_\ell$ est un foncteur exact à gauche nul sur les groupes finis), c'est un \deleted{groupe} \uline{module libre sur} $\mathbb Z_\ell$ de rang égal : $2\alpha+\deleted{\unsure{2}}\mu$, où $\alpha$ ($\mu$) est la dim. de la ``partie'' ``abélienne'' resp. \inserted{``partie''} ``torique'' de $(A_K)^\circ_{\mathrm{red}}$. Ici encore, $\pi$ opère sur $E$ par transport de structure.

\uline{Remarque} 0.5. On peut regarder l'exemple $0.\unsure{3}$ comme étant essentiellement \deleted{un} \inserted{le} cas particulier de l'exemple 0.4, \inserted{pour $i=1$}.
De façon précise, \inserted{on montre que} on a un isom canonique
\[
 H^1(A_{\bar K},\mathbb Z_\ell)\simeq \deleted{\unsure{H^1}}\operatorname{Hom}_{\mathbb Z_\ell}(T_\ell(A(\bar K)),\mathbb Z_\ell),
\]
ou ce qui revient au même, puisque $T_\ell(A_{\bar K})$ est libre de type fini sur $\mathbb Z_\ell$ :
\[
 T_\ell(A(\bar K))=\operatorname{Hom}_{\mathbb Z_\ell}(H^1\unsure{(}A_{\bar K},\mathbb Z_\ell),\mathbb Z_\ell).
\]
D'autre part, \deleted{dans 0.3, si} si \deleted{$i=1$} dans 0.3. $X$ est propre et $i=1$, on trouve
\[
\begin{aligned}
 H^1(X_{\bar K},\mathbb Z_\ell(1))&\simeq\varprojlim_\nu H^1(X_{\bar K},\mu_{\ell^\nu})\simeq T_\ell(\operatorname{Pic}(X_{\bar K}))\\
 &\simeq T_\ell(\underline{\operatorname{Pic}}_{X/K}(\bar K))
\end{aligned}
\]
\margin{3}
